\begin{frame}{Popper의 반증주의: 왜 ``반박 가능''인가}
  \begin{columns}[T]
    \begin{column}{0.55\textwidth}
      Karl Popper, 1934년 \textit{The Logic of Scientific
      Discovery}:
      \begin{block}{반증주의 (falsifiability)}
        한 주장은 ``그것이 틀렸음을 보여주는 관측이
        \textbf{원칙적으로} 상상이 가능한가''로 판별된다.
        \textbf{반박될 수 없는 주장은 과학이 아니다.}
      \end{block}
    \end{column}
    \begin{column}{0.43\textwidth}
      \kb{peer-review의 적용}
      좋은 질문이란 \textbf{발표팀이 제대로 답하지 못하면
      결과 자체가 흔들리는} 질문:
      \begin{itemize}
        \item ``더 잘 설명해 주세요'' = \textbf{요청}.
        \item ``Week 6 교차검증 없이 $\lambda$를 정했는데, 그
          검증 성능이 우연일 가능성은 없는가?'' = \textbf{반박}.
      \end{itemize}
    \end{column}
  \end{columns}
\end{frame}
